\selectlanguage{ngerman}%
\frenchspacing%
\begin{ledgroupsized}[r]{120mm}%
\footnotesize%
\pstart%
\noindent\textbf{Überlieferung:}%
\pend%
\end{ledgroupsized}%
\begin{ledgroupsized}[r]{114mm}%
\footnotesize%
\pstart%
\parindent -6mm%
\makebox[6mm][l]{\textit{L}}%
Aufzeichnung:
LH XXXVII 4, Bl.\;171.
Ein Blatt 4\textsuperscript{o},
wiederverwendeter Umschlag eines nicht ermittelten Briefes;
\href{\Scan{LH_37_04_171wz}}{Wassezeichen}. %
Eine Seite auf Bl.\;\href{\Scan{LH_37_04_171r}}{171\,r\textsuperscript{o}};
der letzte Absatz (\refpassage{LH_37_04_171_171r_job-1}{LH_37_04_171_171r_job-2}) weist einen verschiedenen Schreibduktus auf.
Bl.\;\href{\Scan{LH_37_04_171v}}{171\,v\textsuperscript{o}} überliefert folgende Briefanschrift von fremder Hand:
\textit{A Monsieur}
\lbrack/\rbrack\
\textit{Monsieur Leibniz Conseileur}
\lbrack/\rbrack\
\textit{de la Cour de S.\;A.\;El.\;de Br.}
\lbrack/\rbrack\
\textit{et Lunebourg p}
\lbrack/\rbrack\
\textit{à Hannover.}%
\indexO{Braunschweig}\indexO{Luneburg}\indexO{Hannover}
\pend%
\end{ledgroupsized}%
\vspace{5mm}%
\begin{ledgroup}%
\footnotesize%
\pstart%
\noindent%
\textbf{Datierungsgründe:}
Die Datierung der Aufzeichnung N.\;\ref{53451} beruht auf äußeren Eigenschaften des Text\-trägers %
sowie auf inhaltlichen Betrachtungen.
Siehe die ausführliche Begründung auf S.~\pageref{53451_dat}.%
\pend%
\end{ledgroup}%
\selectlanguage{latin}%
\frenchspacing%
\vspace{8mm}%
\normalsize%
\pstart%
\noindent%
\lbrack%
\href{\Scan{LH_37_04_171r}}%
{171\,r\textsuperscript{o}%
}%
\rbrack\ %
Dn.
\edtext{%
Amontons%
\indexN{\textls{Amontons}, Guillaume 1663\textendash1705}
in \textit{Ephemerid.\;Paris.}%
\indexS{ephemerides}%
\;\edtext{%
10 Maji%
}{%
\lemma{\textit{Paris.}\;10}%
\Bfootnote{%
\textit{(1)}~Mart. %
\textit{(2)}~Maji%
~\textit{L}%
}}
1688
et \textit{Act.\;erud.}
inde desumtis Julii 1688%
}{%
\lemma{Amontons \lbrack...\rbrack\ 1688}%
\Cfootnote{%
Siehe
\textsc{G.\;Amontons},
\glqq Extrait d'une lettre touchant la construction d'un nouveau tube pour faire le vuide à une si petite hauteur perpendiculaire qu'on voudra\grqq\
(\textit{JS}, 10.\;Mai 1688, S.\;394--396)\cite{01511}
sowie die anonyme Übersetzung des Auszugs:
\glqq Contenta epistolae de constructione novi tubi ad faciendum vacuum in altitudine perpendiculari tam exigua quam placuerit\grqq\
(\textit{AE} 1688, S.\;376--378 u.\;Tafel VII, Abb.\;3; Übersetzer war C.\;Pincker).%
\indexN{\textls{Pincker} (Pinckert), Christoph gest. 1692}%
}}
proposuit modum exhibendi vacuum%
\indexS{modus exhibendi vacuum}%
\indexS{vacuum exhibendum}%
\indexS{vacuum factum in tubo}
altitudine liquoris dati,%
\indexS{liquor in tubo}
quantumlibet parva.%
\indexS{altitudo liquoris in tubo}
Nempe loco altitudinis continuae unius tubi%
\indexS{tubus plenus liquore}
adhibet
plures tubos
\edtext{%
aequales%
}{%
\lemma{aequales}%
\Bfootnote{%
\textit{erg.\;L}%
}}
liquore plenos
simul altitudinem requisitam
componentes;
inter se connexos ope tuborum
solo aere plenorum;%
\indexS{tubus plenus aere}%
\indexS{aer replens tubum}
sed elastri%
\indexS{elastrum aeris}
quod una cum
\edtext{%
altitudine liquoris%
\indexS{altitudo liquoris in tubo}%
\lbrack,\rbrack%
}{%
\lemma{altitudine}%
\Bfootnote{%
\textit{(1)}~\mercury ii %
\textit{(2)}~liquoris%
~\textit{L}%
}}
aeri externo incumbenti%
\indexS{aer externus incumbens}
resistat.%
\indexS{elastrum resistens aeri incumbenti}
\pend%
\pstart%
\edtext{%
Ponamus
nos uti velle Mercurio%
\indexS{mercurius in tubo}
et quidem altitudine 14 pollicum;%
\indexS{pollex mercurii}
dividamus 28
ordinariam altitudinem \mercury ii,
\edtext{%
per 14,
quotiens est 2.%
\indexS{quotiens}%
}{%
\lemma{per~14,}%
\Bfootnote{%
\textit{(1)}~fit~2 %
\textit{(2)}~quotiens est~2.%
~\textit{L}%
}}
Ergo duobus Tubis%
\indexS{tubus plenus aere}
est opus,
unoquoque pollicum 14.
\edtext{%
Nempe \textit{BC}%
}{%
\lemma{Nempe \textit{BC}}%
\Cfootnote{%
Vgl.\;das Diagramm \lbrack\textit{Fig.\:1}\rbrack.%
}}
altitudo 14 pollicum,
tubus \textit{BC} plenus Mercurio,%
\indexS{mercurius in tubo}%
\indexS{tubus plenus mercurio}
ut et similis tubus \textit{DE},
altitudinis similiter 14 pollicum;%
\indexS{altitudo mercurii in tubo}
qui connexi Tubo \textit{CD} pleno aere,%
\indexS{tubus plenus aere}
cujus elastrum%
\indexS{elastrum aeris}
una cum pondere Mercurii 14 pollicum%
\indexS{pondus mercurii in tubo}%
\indexS{pollex mercurii}
sustineat atmosphaerae pondus.%
\indexS{pondus atmosphaerae}
Porro tubus compositus \textit{ABCDEF}%
\indexS{tubus compositus} %
sit apertus in \textit{A},
clausus in \textit{F},
fiet vacuum in \textit{FE},%
\indexS{vacuum factum in tubo}
perinde ac si \textit{BC} et \textit{DE}
in unam altitudinem essent conjuncta.
Si 7 pollicum altitudinis%
\indexS{altitudo mercurii in tubo}%
\indexS{pollex mercurii}
deberet esse \textit{BC},
vel \textit{DE},
quatuor Tubis opus foret.%
\indexS{tubus plenus mercurio}%
}{%
\lemma{Ponamus \lbrack...\rbrack\ foret}%
\Cfootnote{%
Vgl.
\textsc{Amontons},
\glqq Extrait d'une lettre\grqq\
(S.\;394--396);\cite{01511}
\glqq Contenta epistolae\grqq\
(S.\;377\,f.).%
}}%
\pend%
\newpage%
\vspace*{0.0em}%
\begin{edfigure}[0.0em][2.00em]%
\includegraphics[width=0.38\textwidth]%
{gesamttex/edit_VIII-4/images/LH_37_04_171_d_171r.pdf}%
\edcaption[0.50em]
{\lbrack \textit{Fig.\:1}\rbrack%
\edtext{}{%
\lemma{\lbrack \textit{Fig.\:1}\rbrack}%
\Cfootnote[nonum]{%
Vgl. das Diagramm in:
\textsc{Amontons},
\glqq Extrait d'une lettre\grqq\
(S.\;395);\cite{01511}
\glqq Contenta epistolae\grqq\
(Tafel VII, Abb.\;3.)%
}}%
}%
\end{edfigure}%
\pstart%
Sed tunc invenio difficultatem%
\indexS{difficultas de elastro aeris}
in elastro aeris definiendo%
\indexS{elastrum aeris}%
\lbrack;\rbrack\
nam,
ut aequivaleat
\edtext{%
atmosphaerae%
\indexS{atmosphaera externa}
cum 7 pollicibus \mercury ii,%
\indexS{mercurius in tubo}%
\indexS{pollex mercurii}
debet per \textlangle se\textrangle\ 
aequivalere%
}{%
\lemma{atmosphaerae}%
\Bfootnote{%
\textit{(1)}~debet tale esse, ut %
\textit{(a)}~\textlangle cum\textrangle\ %
\textit{(b)}~cum %
\textit{(2)}~cum 7 %
\lbrack...\rbrack\ \textlangle se\textrangle\ aequivalere%
~\textit{L}%
}}
21 pollicibus Mercurii%
\indexS{pollex mercurii}%
\lbrack;\rbrack\
ita sane ab atmosphaera externa%
\indexS{atmosphaera externa}
non patietur.
Verum ita non sustinebit tantum,
sed et attollet Mercurium tubi ultimi \textit{ED},%
\indexS{mercurius in tubo}
adeoque nec eum descendere%
\indexS{descensus mercurii in tubo}
seu vacuum facere patietur.%
\indexS{vacuum factum in tubo}
Itaque vereor
ne auctoris inventum succedat.%
\indexS{auctor}
Solummodo in tubis duobus,%
\indexS{tubus compositus}
seu 14 pollicibus;
ita enim aer \textit{CD}
simul et aeri externo%
\indexS{aer externus}
cum 14 pollicibus aequivalet,
et ipsis 14 pollicibus.%
\indexS{pollex mercurii}
\pend%
\pstart%
Caeterum utcunque res sit,
non sufficit ista methodus%
\indexS{methodus insufficiens}
ad barometrum%
\indexS{barometrum brevius}
brevius solito efficiendum,
quia aer inclusus in \textit{CD}%
\indexS{aer inclusus tubo}
a calore et frigore nimium pateretur.%
\indexS{calor}%
\indexS{frigus}
Quare aliud remedium cogitavi,%
\indexS{remedium cogitatum}
quo et caloris frigorisque susceptibilitas,%
\indexS{susceptibilitas caloris}%
\indexS{susceptibilitas frigoris}%
\indexS{calor}%
\indexS{frigus}
et ipsa illa difficultas de aeris Elastro%
\indexS{difficultas de elastro aeris}
amovetur.
\pend%
\newpage%
\pstart%
Nempe ponamus exempli causa
desiderari Barometrum altitudinis pollicum 4.%
\indexS{barometrum brevius}%
\indexS{altitudo barometri}
\edtext{%
Erit illud tuborum%
}{%
\lemma{Erit}%
\Bfootnote{%
\textit{(1)}~illud pollicum %
\textit{(2)}~illud tuborum%
~\textit{L}%
}}
septem per sexalios connexorum
seu tuborum 13.%
\indexS{tubus compositus}
Porro \textit{CD} tubum connectentem%
\indexS{tubus connectens tubos mercuriales}
\mercury iales%
\indexS{tubus mercurialis}
implebimus non aere,%
\indexS{aer implens tubum}
sed liquore%
\indexS{liquor implens tubum}
congelationis%
\indexS{congelatio}%
\indexS{liquor expers congelationis}
et compressionis%
\indexS{compressio}%
\indexS{liquor expers compressionis}
ac rarefactionis%
\indexS{rarefactio}%
\indexS{liquor expers rarefactionis}
\edtext{%
expert[e].%
}{%
\lemma{expertae}%
\Bfootnote{%
\textit{L\;ändert Hrsg.}%
}}
Tale Barometrum%
\indexS{barometrum brevius}
poterit exiguae pyxidi includi,%
\indexS{pyxis in barometro}
praesertim si absit ventricositas.%
\indexS{ventricositas}
\edlabel{LH_37_04_171_171r_zfrw-1}%
In fine et initio tuborum%
\indexS{tubus compositus}
ventricositate retenta%
\indexS{ventricositas}
possemus super \textit{A} collocare liquorem%
\indexS{liquor implens tubum}
\edtext{%
methodo Hugeniana;%
\indexN{\textls{Huygens} (Hugenius, Ugenius, Hugens, Huguens), Christiaan 1629\textendash1695}%
\indexS{methodus Hugeniana}%
}{%
\lemma{methodo Hugeniana}%
\Cfootnote{%
Siehe
C.\;\textsc{Huygens},
\glqq Extrait d'une autre lettre touchant une nouvelle maniere de Barometre\grqq\
(\textit{JS}, 12.~Dezember 1672, S.\;153\,f.;%
\cite{01522}
\textit{HO} VII, S.\;239).%
\cite{00113}%
}}
sed sublata illa%
\lbrack,\rbrack\
quod praestat,
licebit
\edlabel{LH_37_04_171_171r_hooke-1}
\edtext{%
Hookiana%
\indexN{\textls{Hooke} (Hookius, Hook), Robert 1635\textendash1703}%
\indexS{methodus Hookiana}%
}{%
\lemma{Hookiana}%
\Cfootnote{%
Siehe \textsc{R.\;Hooke},
\textit{Micrographia}, Preface
(London 1665, S.\:[c2]--[c3] u.\;Tafel I, Abb.\;I);%
\cite{00061}
\textsc{Ders.},
\glqq A new Contrivance of Wheel-Barometer\grqq\
(\textit{PT}, I, 4.\;Juni 1666,
S.\;218\,f. u.\;Abb.\;1 vor S.\;215).%
\cite{01523}
Auf Hookes Radbarometer beziehen sich in diesem Band auch die Aufzeichnungen
N.\;\ref{57038} (\refpassage{LH_37_04_157-158_157v_udz-1}{LH_37_04_157-158_157v_udz-2})
und
N.\;\ref{57041} (\refpassage{LH_37_04_162-163_163r_hooke-1}{LH_37_04_162-163_163r_hooke-2} u.\;\refpassageohne{LH_37_04_162-163_163r_hooke-3}{LH_37_04_162-163_163r_hooke-4})
vom Herbst 1677.%
}}
vel simili latiore
uti,
\edtext{%
dum grave innatans%
}{%
\lemma{dum}%
\Bfootnote{%
\textit{(1)}~innatans %
\textit{(a)}~aliquod %
\textit{(b)}~gra %
\textit{(2)}~grave innatans%
~\textit{L}%
}}
Mercurio in \textit{B},%
\indexS{grave innatans mercurio}%
\indexS{mercurius in barometro}
secum
\edtext{%
indicem circumagit,%
\indexS{index barometri}
vel virunculum%
\indexS{virunculus}
attollit aut deprimit.%
}{%
\lemma{indicem}%
\Bfootnote{%
\textit{(1)}~vel virunculum mo %
\textit{(2)}~circumagit vel %
\lbrack...\rbrack\ aut deprimit.%
~\textit{L}%
}}%
\edlabel{LH_37_04_171_171r_hooke-2}%
\edlabel{LH_37_04_171_171r_zfrw-2}%
\pend%
\pstart%
\edtext{%
Liquor quo Hugenius%
\indexN{\textls{Huygens} (Hugenius, Ugenius, Hugens, Huguens), Christiaan 1629\textendash1695}
utitur in suo Barometro,%
\indexS{barometrum Hugenii}%
\indexS{liquor in barometro}%
}{%
\lemma{Liquor \lbrack...\rbrack\ Barometro}%
\Cfootnote{%
Siehe
\textsc{Huygens},
\glqq Extrait d'une autre lettre\grqq\
(S.\;154;%
\cite{01522}
\textit{HO} VII, S.\;239).%
\cite{00113}%
}}
etiam nobis utlilis
\edlabel{LH_37_04_171_171r_aud-1}%
erit.%
\edtext{}{%
\xxref{LH_37_04_171_171r_aud-1}{LH_37_04_171_171r_aud-2}%
\lemma{erit.}%
\Bfootnote{%
\lbrack/\rbrack\ %
\textit{(1)}~Pour d %
\textit{(2)}~Mons.\;Lami %
\textit{(3)}~Le pere Lami%
~\textit{L}%
}}%
\pend%
\vspace{1.0em}%
\pstart%
\noindent%
\lbrack\textit{Nachfolgender Absatz in L möglicherweise nachträglich ergänzt:}\rbrack\
\pend%
\vspace{0.5em}%
\pstart%
\noindent%
\edlabel{LH_37_04_171_171r_job-1}%
\edtext{%
Le pere Lami%
\indexN{\textls{Lamy} (L'Amy, Lami), Bernard (Bernhard) 1640\textendash1715}%
\edlabel{LH_37_04_171_171r_aud-2}
(chez M.\;Vallemont%
\indexN{\textls{Le Lorrain de Vallemont} (Vallemontius), Pierre 1649\textendash1721}
\textit{tr.\;de la baguette} p.\;262)
\edtext{%
dit que la liqueur
qu'on%
}{%
\lemma{dit}%
\Bfootnote{%
\textit{(1)}~qu'on %
\textit{(2)}~que la liqueur qu'on%
~\textit{L}%
}}
peut employer pour les Barometres%
\indexS{liquer pour le barometre}
\edtext{%
doubles,%
\indexS{barometre double}
et \textit{qui}%
}{%
\lemma{doubles,}%
\Bfootnote{%
\textit{(1)}~qui %
\textit{(2)}~et \textit{qui}%
~\textit{L}%
}}
\textit{ne gele point en hyver%
\indexS{liquer qui non gele pas}%
\indexS{hyver}
et ne puisse dissoudre le vif argent},%
\indexS{argent vif}
pourra estre \textit{l'eau fort%
\indexS{eau fort}
melée avec six fois autant d'eau commune}.%
\indexS{eau commune}%
}{%
\lemma{Le pere [...] commune}%
\Cfootnote{%
Siehe
\textsc{P.\;Le Lorrain de Vallemont},
\textit{Physique occulte}, %
ch.\;XI (Amsterdam 1693, S.\;262);%
\cite{01513}
vgl.
\textsc{B.\;Lamy}, \textit{Traitez de mechanique},
traité II, prop.\;6, probl.\;1; prop.\;7, probl.\;2
(Paris 1679, S.\;196--198, 216--219).\cite{01476}
Die von Leibniz getreu wiedergegebenen Worte
\textit{qui ne gele} \lbrack...\rbrack\ \textit{d'eau commune}
sind nahezu unverändert aus
\textsc{Huygens},
\glqq Extrait d'une autre lettre\grqq\
(S.\;154;%
\cite{01522}
\textit{HO} VII, S.\;239)%
\cite{00113}
übernommen.%
}}%
\edlabel{LH_37_04_171_171r_job-2}%
\pend%
\vspace{5.5em}%
\pstart%
\noindent%
\centering%
*     *     *
\pend%
\newpage%
\selectlanguage{german}%
\frenchspacing%
\footnotesize%
\pstart%
\noindent%
\textls{Ausführliche Datierungsgründe:}%
\pagelabel{53451_dat}
Die titellose und undatierte Aufzeichnung N.\;\ref{53451}, die nahezu den Charakter eines Exzerpten aufweist, ist dem Modell eines \glqq kürzeren\grqq\ bzw. \glqq zusammengesetzten\grqq\ Barometers gewidmet, welches Guillaume Amontons\indexN{\textls{Amontons}, Guillaume 1663\textendash1705} im Jahre 1688 beschrieben hatte.
Dieses wird im Text kurz geschildert, kritisch erörtert und mit Ideen zu technischen Innovationen, die Christiaan Huygens\indexN{\textls{Huygens} (Hugenius, Ugenius, Hugens, Huguens), Christiaan 1629\textendash1695} und Robert Hooke\indexN{\textls{Hooke} (Hookius, Hook), Robert 1635\textendash1703} Jahre zuvor erarbeitet hatten, erweitert.
Die vorgeschlagene Datierung der Aufzeichnung ergibt sich aus folgenden Beobachtungen:
\pend%
\pstart%
(1)
Das im Textträger der Aufzeichnung N.\;\ref{53451} vorliegende \href{\Scan{LH_37_04_171wz}}{Wassezeichen} ist bislang nicht ermittelt worden.
Bei diesem Textträger handelt es sich allerdings um den wiederverwendeten Umschlag eines Briefes an Leibniz, dessen Absender und Absendeort noch unbekannt sind.
Leibniz wird in der Briefanschrift auf Bl.\;\href{\Scan{LH_37_04_171v}}{171\,v\textsuperscript{o}} als Hofrat \textit{de Son Altesse Electorale de Brunschvick et Lunebourg} genannt.\indexO{Braunschweig}\indexO{Luneburg}
Sein zweiter Hannoveraner\indexO{Hannover} Dienstherr, der Herzog Ernst August von Braunschweig-Lüneburg,\indexN{\textls{Braunschweig-Lüneburg}, Ernst August von, Herzog und Kurfürst von Hannover 1679\textendash1698} erwarb offiziell die Kurwürde gegen Ende Dezember 1692 (\textit{Chronik}, S.\;120).\cite{01236}
Daher kann der Brief an Leibniz, auf dessen Umschlag die Aufzeichnung N.\;\ref{53451} verfasst wurde, nicht vor 1693 abgeschickt worden sein.
Das soll als frühester Terminus post quem der Datierung von N.\;\ref{53451} betrachtet werden.
\pend%
\pstart%
(2)
Kurfürst Georg Ludwig von Braunschweig-Lüneburg,\indexN{\textls{Braunschweig-Lüneburg}, Georg Ludwig von, Kurfürst von Hannover 1698\textendash1727} Ernst Augusts\indexN{\textls{Braunschweig-Lüneburg}, Ernst August von, Herzog und Kurfürst von Hannover 1679\textendash1698} Nachfolger als Dienstherr Leibnizens, stieg im August 1714 auf Englands\indexO{England} Thron und wurde somit in Personalunion König von Großbritannien\indexO{Großbritannien} und Irland.\indexO{Irland}
Da dieser Umstand in der besagten Briefanschrift keine Berücksichtigung findet, ist unwahrscheinlich, dass der Brief nach dem Sommer 1714 gesendet wurde.
Dies dürfte somit als der Terminus ante quem der Datierung der Aufzeichnung N.\;\ref{53451} gelten.
\pend%
\pstart%
(3)
Der Terminus post quem der Datierung lässt sich aber möglicherweise weiter bestimmen.
Es gilt nämlich zu bemerken, dass in N.\;\ref{53451} (\refpassage{LH_37_04_171_171r_zfrw-1}{LH_37_04_171_171r_zfrw-2}) Leibniz auf dieselben Aufsätze von Huygens\indexN{\textls{Huygens} (Hugenius, Ugenius, Hugens, Huguens), Christiaan 1629\textendash1695} (\glqq Extrait d'une autre lettre touchant une nouvelle maniere de Barometre\grqq, \textit{JS}, 12.~Dezember 1672, S.\;152--156;\cite{01522} \textit{HO} VII, S.\;238--241)\cite{00113} und Hooke\indexN{\textls{Hooke} (Hookius, Hook), Robert 1635\textendash1703} (\textit{Micrographia}, Preface, London 1665, S.\:[c3];\cite{00061} \glqq A new Contrivance of Wheel-Barometer\grqq, \textit{PT}, 4. Juni 1666, S.\;218\,f.)\cite{01523} verweist wie in der umfangreicheren Aufzeichnung N.\;\ref{60267} im vorliegenden Band (\refpassage{LH_37_04_149-150_149r_rvb-1}{LH_37_04_149-150_149r_rvb-2}).
Diese letztere lässt sich auf den Zeitraum von Mitte September 1696 bis zur zweiten Hälfte 1697 datieren, wobei der Schlussteil bis zur zweiten Hälte 1699 entstanden sein könnte (siehe hierzu die einleitende Erläuterung zu N.\;\ref{paradoxon_intro}, S.\;\pageref{paradoxonbarometricum_intro}\,ff.). %
Es ist bemerkenswert, dass in der ausführlicheren und an Literaturhinweisen großzügigen Aufzeichnung N.\;\ref{60267} Amontons'\indexN{\textls{Amontons}, Guillaume 1663\textendash1705} Barometer, ungeachtet seiner Verwandschaft mit den von Huygens\indexN{\textls{Huygens} (Hugenius, Ugenius, Hugens, Huguens), Christiaan 1629\textendash1695} und Hooke\indexN{\textls{Hooke} (Hookius, Hook), Robert 1635\textendash1703} beschriebenen Vorrichtungen, nicht erwähnt wird.
Dieser Umstand ist Grund für die Vermutung, dass Leibniz zu dem Zeitpunkt, als er N.\;\ref{60267} verfasste, sich mit Amontons'\indexN{\textls{Amontons}, Guillaume 1663\textendash1705} Bericht im \textit{Journal des sçavans} bzw. den \textit{Acta eruditorum} noch nicht gründlich befasst hatte.
Demgemäß dürfte auch die vorliegende Aufzeichnung N.\;\ref{53451} erst nach N.\;\ref{60267} entstanden sein.
\pend%
\pstart%
(4)
Zum gleichen Schluss kommt man auch, wenn man den expliziten Verweis auf Pierre Le Lorrain de Vallemonts\indexN{\textls{Le Lorrain de Vallemont} (Vallemontius), Pierre 1649\textendash1721} Abhandlung \textit{La physique occulte ou traité de la baguette divinatoire}\cite{01513} in Betracht zieht, welcher in dem möglicherweise nachträglich hinzugefügten Schlussabsatz der Aufzeichnung N.\;\ref{53451} (\refpassage{LH_37_04_171_171r_job-1}{LH_37_04_171_171r_job-2}) anzutreffen ist.
Vallemonts\indexN{\textls{Le Lorrain de Vallemont} (Vallemontius), Pierre 1649\textendash1721} Abhandlung war nicht vor dem 21.\;März 1693 in Paris\indexO{Paris} erschienen, wie man dem Kolophon der Erstausgabe entnehmen kann.
Wie die Seitenangabe in dem besagten Verweis zeigt, hatte Leibniz offenbar die zweite, im selben Jahr in Amsterdam\indexO{Amsterdam} gedruckte Ausgabe der \textit{Physique occulte}\cite{01513} vor Augen, die wohl nicht vor der zweiten Jahreshälfte erschienen war.
Freilich war Leibniz über Vallemonts\indexN{\textls{Le Lorrain de Vallemont} (Vallemontius), Pierre 1649\textendash1721} Abhandlung bereits kurz nach deren Erscheinen in Paris\indexO{Paris} durch Claude Nicaises\indexN{\textls{Nicaise} (Nicasius), Claude 1623\textendash1701} Brief vom 23.\;April 1693 informiert worden (\textit{LSB} II,\;2 N.\;216, S.\;690.23--26)\cite{02518} und verweist auf sie zum ersten Mal -- ohne den Titel ausdrücklich zu nennen -- in seinem Schreiben an Wilhelm Ernst Tentzel\indexN{\textls{Tentzel} (Tenzelius), Wilhelm Ernst 1659\textendash1707} vom 5./15. Mai 1693 (\textit{LSB} I,\;9 N.\;266, S.\;424.12--15).\cite{01467}
Keine Erwähnungen der \textit{Physique occulte}\cite{01513} im Briefwechsel zeigen jedoch mit Evidenz, dass Leibniz direkte Kenntnis von ihr hatte; vgl. das Schreiben an Wilhelm Ernst Tentzel\indexN{\textls{Tentzel} (Tenzelius), Wilhelm Ernst 1659\textendash1707} vom 16./26. Januar 1694 (\textit{LSB} I\;10 N.\;132, S. 226.21--24)\cite{01468}
und dasjenige an Henri Basnage de Bauval\indexN{\textls{Basnage de Bauval} (Bellavallius, Beauval, Bauvallius), Henri 1656\textendash1710} von Februar/März 1694 (\textit{LSB} II,\;2 N.\;258, S.\;794.7--14).\cite{01469}
Ausschlaggebend für den Verweis auf die \textit{Physique occulte}\cite{01513} in der Aufzeichnung N.\;\ref{53451} könnte vielmehr der Umstand gewesen sein, dass Bernardino Ramazzini\indexN{\textls{Ramazzini} (Ramazzinus), Bernardino 1633\textendash1714} sie in seinen 1695 in Modena\indexO{Modena}\cite{01516} und 1696 in Augsburg\indexO{Augsburg}\cite{01517} veröffentlichten \textit{Ephemerides barometricae} nennt, von deren Lektüre wiederum Leibnizens Aufzeichnung N.\;\ref{60267} herrührt (siehe erneut die einleitende Erläuterung zu N.\;\ref{paradoxon_intro}, S.\;\pageref{paradoxonbarometricum_intro}\,ff.).
Ramazzini\indexN{\textls{Ramazzini} (Ramazzinus), Bernardino 1633\textendash1714} erwähnt dort (Modena 1695, S.\;38\cite{01516} bzw. Augsburg 1696, S.\;15)\cite{01517} ausdrücklich dasselbe Kapitel~XI von Vallemonts\indexN{\textls{Le Lorrain de Vallemont} (Vallemontius), Pierre 1649\textendash1721} Abhandlung, wie Leibniz im besagten Schlussabsatz der vorliegenden Aufzeichnung N.\;\ref{53451}.
Demnach ist nicht unwahrscheinlich, dass Leibniz erst im Zuge seiner Auseinandersetzung mit Ramazzinis\indexN{\textls{Ramazzini} (Ramazzinus), Bernardino 1633\textendash1714} \textit{Ephemerides barometricae}\cite{01516}\cite{01517} das dem Barometer gewidmete Kapitel~XI der \textit{Physique occulte}\cite{01513} las.
Auch diese Betrachtung ist Grund für die Vermutung, dass N.\;\ref{53451} nicht vor der Aufzeichnung N.\;\ref{60267} und somit nicht vor der zweiten Hälfte September 1696 entstand.
\pend%
\pstart%
Aus den angeführten Gründen ergibt sich als vorgeschlagene Gesamdatierung der vorliegenden Aufzeichnung N.\;\ref{53451} die Zeitspanne von der zweiten Hälfte September 1696 bis zum Sommer 1714.
Eine frühere Entstehung von N.\;\ref{53451} ist zwar möglich, aber erst ab 1693;
wenigstens der Schlussabsatz der Aufzeichnung (\refpassage{LH_37_04_171_171r_job-1}{LH_37_04_171_171r_job-2}) muss nach Leibnizens Lektüre von Vallemonts\indexN{\textls{Le Lorrain de Vallemont} (Vallemontius), Pierre 1649\textendash1721} \textit{Physique occulte}\cite{01513} verfasst worden sein, d.\,h. frühestens in der zweiten Hälfte 1693.

\pend%
\selectlanguage{latin}%
\frenchspacing%
\normalsize%
\newpage%
